\section{Objective}

This experiment uses the \texttt{aloeL.jpg} and \texttt{aloeR.jpg} images from the official OpenCV sample data to implement a complete pipeline from a stereo image pair to a disparity map and a relative-depth map. The focus is not on reimplementing a stereo-matching algorithm from scratch, but on connecting input validation, preprocessing, matching, numerical filtering, result saving, and evidence checking into a reproducible experimental loop.

\subsection{Course relevance}

Stereo depth estimation connects the geometric models of computer vision with practical image-processing interfaces. Disparity is the directly observable pixel displacement between two images, whereas depth is also affected by the focal length, baseline, and calibration quality. This experiment distinguishes between producing an image with visually layered structure and obtaining depth with physical units that can be used for measurement. It also illustrates how invalid matches, occlusions, and texture conditions affect the result.

\subsection{Experimental objectives}

\begin{enumerate}
  \item Understand the pinhole stereo model, the epipolar constraint, and the relationship between disparity and depth;
  \item Learn the OpenCV interfaces for image loading, grayscale preprocessing, and \texttt{StereoSGBM} stereo matching;
  \item Convert the fixed-point disparity returned by the matcher into floating-point disparity in pixel units, and perform valid-region filtering and visualization;
  \item Explain how parameters such as \texttt{minDisparity}, \texttt{numDisparities}, the smoothness penalties, and speckle filtering affect the matching result;
  \item When camera calibration parameters paired with the image pair are unavailable, compute and analyze a relative-depth proxy without physical units;
  \item Save the left and right images, disparity map, valid-disparity mask, arrays, and parameter summary generated by the real run, and use them to verify the results.
\end{enumerate}

\subsection{Reproducibility and input/output scope}

This report uses the aloe stereo image pair from OpenCV's official \texttt{4.x/samples/data} directory as the only primary experimental input. The program command, software versions, matching parameters, output-array dimensions, and statistics are taken from \codepath{task2/results/params.json} and \codepath{run_evidence.md} generated by this run. The outputs include pixel disparity and \(\mathrm{depth\_proxy}=1/d\) on an arbitrary scale. Because the current image pair has no verified focal length, baseline, or $\mathbf{Q}$ matrix, this proxy is not interpreted as metric depth in metres, centimetres, or any other physical unit.
